\section*{Chapter \ref{ch:b1:precipitation} --- Precipitation and Solubility}

\begin{solution}{exo:b1:precipitation:1}
\ce{BaSO4(s) <=> Ba^2+ + SO4^2-}, $K_s = [\ce{Ba^2+}][\ce{SO4^2-}] =
10^{-9.97}$. \ce{CaF2(s) <=> Ca^2+ + 2F-}, $K_s = [\ce{Ca^2+}][\ce{F-}]^2 =
10^{-9.84}$. \ce{Ag2CrO4(s) <=> 2Ag+ + CrO4^2-}, $K_s =
[\ce{Ag+}]^2[\ce{CrO4^2-}] = 10^{-11.95}$. \ce{Fe(OH)3(s) <=> Fe^3+ + 3OH-},
$K_s = [\ce{Fe^3+}][\ce{OH-}]^3 = 10^{-38.55}$. For $K_s = \num{2.0e-13}$,
$\mathrm{p}K_s = 12.70$.
\end{solution}

\begin{solution}{exo:b1:precipitation:2}
$s = \sqrt{K_s} = 10^{-12.27/2} = \qty{7.3e-7}{mol/L}$, that is
$\num{7.3e-7} \times 187.8 = \qty{1.4e-4}{g/L} = \qty{0.14}{mg/L}$.
\end{solution}

\begin{solution}{exo:b1:precipitation:3}
After mixing (volume doubled), $[\ce{Ba^2+}] = \qty{1.0e-3}{mol/L}$ and
$[\ce{SO4^2-}] = \qty{1.0e-4}{mol/L}$: $Q_s = \num{1.0e-7} > K_s =
\num{1.1e-10}$, barium sulfate precipitates. In the second case
$[\ce{SO4^2-}] = \num{1.0e-7}$ and $Q_s = \num{1.0e-10} < K_s$: no
precipitate, though only just.
\end{solution}

\begin{solution}{exo:b1:precipitation:4}
In water $s = 10^{-9.97/2} = \qty{1.0e-5}{mol/L}$. In \qty{0.010}{mol/L}
sulfate (\cref{prop:b1:precipitation:common-ion}), $s = 10^{-9.97}/0.010 =
\qty{1.1e-8}{mol/L}$, about a thousand times less (factor 970).
\end{solution}

\begin{solution}{exo:b1:precipitation:5}
In water $s = (10^{-9.84}/4)^{1/3} = \qty{3.3e-4}{mol/L}$, or
\qty{26}{mg/L}. In \qty{0.10}{mol/L} calcium chloride,
$[\ce{Ca^2+}] \approx 0.10$ and $[\ce{F-}] = 2s$, so
$s = \frac12\sqrt{10^{-9.84}/0.10} = \qty{1.9e-5}{mol/L}$. The
approximation $s \ll 0.10$ holds by four orders of magnitude.
\end{solution}

\begin{solution}{exo:b1:precipitation:6}
Onset: $\mathrm{pH} = 14.00 - \frac12(11.25 + \log 0.050) = 14.00 - 4.97 =
9.03$. At \qty{99}{\percent}, $[\ce{Mg^2+}] = \num{5.0e-4}$ and
$\mathrm{pH} = 14.00 - \frac12(11.25 - 3.30) = 10.03$. On a pH axis,
\ce{Mg(OH)2} exists above 9.03; below it, magnesium is entirely \ce{Mg^2+}
at \qty{0.050}{mol/L}.
\end{solution}

\begin{solution}{exo:b1:precipitation:7}
Silver bromide appears first, at $[\ce{Ag+}] = 10^{-12.27 + 2} =
10^{-10.27}$, silver chloride at $10^{-7.75}$. When silver chloride
appears, $[\ce{Br-}] = 10^{-12.27}/10^{-7.75} = 10^{-4.52} =
\qty{3.0e-5}{mol/L}$, \qty{0.30}{\percent} of the bromide. With the usual
criterion of \qty{0.1}{\percent} the separation is not acceptable: the two
$\mathrm{p}K_s$ differ by only 2.5.
\end{solution}

\begin{solution}{exo:b1:precipitation:8}
Silver chromate forms when $[\ce{Ag+}]^2 \times \num{5.0e-3} = 10^{-11.95}$,
$[\ce{Ag+}] = \qty{1.5e-5}{mol/L}$. Then $[\ce{Cl-}] = 10^{-9.75}/\num{1.5e-5}
= \qty{1.2e-5}{mol/L}$, a fraction \num{2.4e-4} (\qty{0.024}{\percent}) of the
chloride: the red colour appears when the chloride is practically all
precipitated.
\end{solution}

\begin{solution}{exo:b1:precipitation:9}
$[\ce{CO3^2-}] = [\ce{HCO3-}]K_{a2}/[\ce{H3O+}] = [\ce{HCO3-}]\,10^{\mathrm{pH}
- \mathrm{p}K_{a2}}$, so $K_s = [\ce{Ca^2+}][\ce{HCO3-}]\,10^{\mathrm{pH} -
\mathrm{p}K_{a2}}$, which is the relation asked. With $[\ce{Ca^2+}] =
[\ce{HCO3-}] = s$: $s^2 = 10^{-8.30 + 10.33 - 8.30} = 10^{-6.27}$,
$s = \qty{7.3e-4}{mol/L}$ (\qty{73}{mg/L} of \ce{CaCO3}), against
$10^{-4.15} = \qty{7.1e-5}{mol/L}$ if the carbonate ion were not a base.
\end{solution}

\begin{solution}{exo:b1:precipitation:10}
1.~Adding \ce{CO2(g) <=> CO2(aq)}: $K' = K K_H = 10^{-4.34 - 1.47} =
10^{-5.81}$.
2.~$[\ce{Ca^2+}] = s$ and $[\ce{HCO3-}] = 2s$, so $4s^3 = K'p$. At
\qty{0.010}{bar}: $s = (10^{-5.81} \times 0.010/4)^{1/3} =
\qty{1.6e-3}{mol/L}$, \qty{157}{mg/L}. Under outdoor air:
$s = \qty{5.5e-4}{mol/L}$, \qty{55}{mg/L}. (Check: under soil air
$[\ce{CO2(aq)}] = 10^{-3.47}$ and $\mathrm{pH} = 6.37 + \log(3.14\times
10^{-3}/3.39 \times 10^{-4}) = 7.34$: hydrogencarbonate does predominate.)
3.~Water saturated in calcite under the high carbon dioxide pressure of
the soil reaches the cave ceiling, where the pressure is lower: it loses
carbon dioxide, $Q > K'$, and calcite precipitates where the drop hangs,
ring after ring.
% ledger: air:CO2, dfg:CO2_g, dfg:CO2_ao (log KH computed in figdata/bachelor-1/aqueous-constants.py)
\end{solution}

\begin{solution}{exo:b1:precipitation:11}
1.~Onset at $\mathrm{pH} = 14.00 - \frac12(16.38 - 2.00) = 6.81$. Complete
redissolution when $[\ce{Zn(OH)4^2-}] = 10^{-2} = 10^{-1.93}[\ce{OH-}]^2$,
$[\ce{OH-}] = 10^{-0.035}$, $\mathrm{pH} = 13.97$. The hydroxide exists
between pH 6.81 and 13.97.
2.~$\mathrm{pH} = 14.00 - 14.45/4 = 10.39$, $s_{\min} = 2 \times
10^{-(16.38 + 1.93)/2} = \qty{1.4e-9}{mol/L}$ (in this two-species model).
3.~Line of slope $-2$ from $(6.81, -2)$ down to the minimum near
$(10.39, -8.9)$, then slope $+2$ up to $(13.97, -2)$; the real curve is
rounded off by the other hydroxo species
(figure in \cref{sec:b1:precipitation:amphoteric}).
\end{solution}

\begin{solution}{exo:b1:precipitation:12}
1.~$\ce{Al(OH)3(s) + OH- <=> Al(OH)4-}$, $K = 10^{-1.23}$; all the
aluminium is dissolved when $10^{-3} \leqslant K[\ce{OH-}]$, that is
$[\ce{OH-}] \geqslant 10^{-1.77}$, $\mathrm{pH} \geqslant 12.23$.
2.~$[\ce{Fe^3+}] = 10^{-38.55}/10^{-3 \times 1.77} = 10^{-33.2}$: the iron
stays entirely as \ce{Fe(OH)3}; filtration separates the iron (solid)
from the aluminium (filtrate).
3.~Magnesium hydroxide is not amphoteric: in excess hydroxide it stays
solid, together with iron(III) hydroxide, and nothing separates.
\end{solution}

\begin{solution}{pb:b1:precipitation:1}
\textbf{1.} \ce{Fe(OH)3(s) <=> Fe^3+ + 3OH-}, $K_s = [\ce{Fe^3+}][\ce{OH-}]^3$;
\ce{Zn(OH)2(s) <=> Zn^2+ + 2OH-}, $K_s = [\ce{Zn^2+}][\ce{OH-}]^2$;
\ce{Mg(OH)2(s) <=> Mg^2+ + 2OH-}, likewise.
\textbf{2.} The constants do not have the same form (power 3 or 2 of
$[\ce{OH-}]$) and the three ions are at different concentrations; only
the onset pH can be compared.
\textbf{3.} At the onset $c[\ce{OH-}]^n = K_s$, so $n\log[\ce{OH-}] =
-\mathrm{p}K_s - \log c$ and $\mathrm{pH} = \mathrm{p}K_e +
\log[\ce{OH-}] = \mathrm{p}K_e - \frac1n(\mathrm{p}K_s + \log c)$.
\textbf{4.} $14.00 - \frac13(38.55 - 3.00) = 2.15$.
\textbf{5.} $14.00 - \frac12(16.38 - 2.30) = 6.96$.
\textbf{6.} $14.00 - \frac12(11.25 - 1.70) = 9.22$.
\textbf{7.} Iron(III) (2.15), then zinc (6.96), then magnesium (9.22). At
pH 2.0, $Q_s = 10^{-3} \times 10^{-36} = 10^{-39} < 10^{-38.55}$ for iron,
and the others are further from precipitating: nothing is solid.
\textbf{8.} \qty{1.0e-6}{mol/L}.
\textbf{9.} $14.00 - \frac13(38.55 - 6.00) = 3.15$.
\textbf{10.} $\log[\ce{Fe^3+}] = -38.55 + 3(14.00 - \mathrm{pH}) = 3.45 -
3\,\mathrm{pH}$: slope $-3$.
\textbf{11.} Up to the onset of zinc hydroxide, pH 6.96.
\textbf{12.} From 3.15 to 6.96.
\textbf{13.} \ce{Fe^3+ + OH- <=> FeOH^2+}, $\beta_1 =
[\ce{FeOH^2+}]/([\ce{Fe^3+}][\ce{OH-}]) = 10^{11.82}$.
\textbf{14.} The two are equal when $[\ce{OH-}] = 10^{-11.82}$, at pH 2.18:
\ce{Fe^3+} predominates below, \ce{FeOH^2+} above.
\textbf{15.} $10^{11.82 + 3.15 - 14.00} = 10^{0.97} = 9.3$: at pH 3.15 the
dissolved iron is 10.3 times the free \ce{Fe^3+}, \qty{1.0e-5}{mol/L}:
only \qty{99.0}{\percent} is removed. The first estimate was not safe.
\textbf{16.} $[\ce{Fe}]_{\mathrm{dissolved}} = 10^{3.45 - 3\,\mathrm{pH}}\,
(1 + 10^{\mathrm{pH} - 2.18})$.
\textbf{17.} $\log[\ce{FeOH^2+}] = 11.82 + \log[\ce{Fe^3+}] + \mathrm{pH} -
14.00 = 1.27 - 2\,\mathrm{pH}$: slope $-2$.
\textbf{18.} $1.27 - 2\,\mathrm{pH} \approx -6.00$ gives 3.64; with the
\ce{Fe^3+} term included, $10^{3.45 - 10.92}(1 + 10^{1.46}) =
\num{1.0e-6}$ at $\mathrm{pH} = 3.64$.
\textbf{19.} $[\ce{Zn^2+}] = \num{5.0e-5}$: $14.00 - \frac12(16.38 - 4.30) =
7.96$. Magnesium hydroxide starts at 9.22: no.
\textbf{20.} \ce{Zn(OH)2(s) + 2OH- <=> Zn(OH)4^2-}, $K = \beta_4 K_s =
10^{14.45 - 16.38} = 10^{-1.93}$.
\textbf{21.} $[\ce{OH-}]^2 = \num{5.0e-3}/10^{-1.93}$, $\mathrm{pH} =
13.81$. The operator must not overshoot with base: above pH 9.2 the
magnesium precipitates with the zinc, and far above, the zinc goes back
into solution.
\textbf{22.} $\num{1.0e-3} \times 1000 \times 0.999 \times 106.8 =
\qty{107}{g}$ per cubic metre.
\textbf{23.} $\num{5.0e-3} \times 1000 \times 0.99 \times 99.4 =
\qty{492}{g}$ per cubic metre.
\textbf{24.} Raised at once above 6.96, both hydroxides precipitate in one
sludge: the zinc is lost in the iron waste, and the waste is
contaminated with zinc. Stopping between the thresholds removes the iron
alone.
\textbf{25.} \textbf{From pH 3.64 to pH 6.96}: below 3.64 more than
\qty{0.1}{\percent} of the iron is still dissolved, as \ce{FeOH^2+} mostly; above
6.96 zinc hydroxide forms. Neglecting \ce{FeOH^2+} would have put the
lower bound at 3.15.
\end{solution}
